\documentclass[10pt]{article}
\usepackage[margin=1in]{geometry}
\usepackage{graphicx} 
\usepackage[latin1]{inputenc}
\usepackage{amsmath,amsfonts,amssymb}
\usepackage{enumitem}
\usepackage{hyperref}
\hypersetup{
	colorlinks=true,
	linkcolor=blue,
	filecolor=magenta,      
	urlcolor=blue,
}

\usepackage{titlesec}
\titlespacing\section{0pt}{12pt plus 4pt minus 2pt}{0pt plus 2pt minus 2pt}


\usepackage{eso-pic}

\setlength\parindent{0pt}
\setlength{\leftskip}{1cm}


\begin{document}


\AddToShipoutPicture*
{\put(350,700){\includegraphics[width=6cm]{OSU_horizontal_2C_O_over_B}}}



\section*{Course Name}
Analytical Workflows

\section*{Course Number}
 IB 516

\section*{Course Credits}
4

\section*{Course Meeting Times}
Lecture, Tuesday \& Thursday 10:00-11:50 am

\section*{Prerequisite and/or Corequisite}
Prerequisite: None

\section*{Course Catalog Description}
Examines and implements the theory and implementation of efficient, reproducible workflows -- 
including best practices in scientific programming, project management, and collaboration --  for 
computational, analytical, and data-driven biological research.

\section*{Course Contents}
\begin{enumerate}[label=\bfseries Week \arabic*:,leftmargin=*,labelindent=3em]
	\itemsep0em
	\item Course overview \& Philosophy\\
	Structuring projects \& Version control with Git
	\item Project proposals\\
	Workflow diagrams
	\item Coding best practices \\
	Hack-a-thon \& discussion
	\item Git w/ GitHub (\emph{Project management \& collaboration})\\
	Hack-a-thon \& discussion
	\item Typesetting with Markdown\\
	Data visualization
	\item Project progress presentations\\
	Hack-a-thon \& discussion
	\item Faster computing (\emph{Vectorization \& parallel computing})\\
	Hack-a-thon \& discussion
	\item Faster computing (\emph{High performance computing})\\
	Hack-a-thon \& discussion
	\item Typesetting with \LaTeX\\\
	Hack-a-thon \& discussion
	\item Project presentations\\
	Project presentations \& Wrap-up
\end{enumerate}

\section*{Course Specific Measurable Student Learning Outcomes}
After successful completion of this course, students will be able to:
\begin{enumerate}[leftmargin=2cm]
	\itemsep0em
	\item Translate a research plan into an explicit analytical workflow;
	\item Apply best practices in scientific programming to construct reproducible research;
	\item Manage and collaborate on complex research projects using a version control system;
	\item Apply analytical workflows to advance their dissertation research.
\end{enumerate}


\section*{Evaluation of Student Performance}
Our primary goal in this course is for students to develop more efficient research skills.  
An important secondary goal is to have students make significant progress on their thesis work.  
Our philosophy is that students can achieve both because our primary goal is best achieved by having 
students practice new tools while working on their own research.
\\\\
Student learning will be verified during weekly ``hack-a-thon'' sessions, by weekly assessment of their 
``pushes'' of ``commits'' to their GitHub repositories, and by evaluation of three project presentations.  
(See rubrics at 
\url{https://github.com/analyticalworkflows/TeachingMaterials/tree/master/course_info/rubrics}).

\begin{itemize}[leftmargin=2cm]
\itemsep0em
\item Hack-a-thon \& Paper discussions \emph{- 200 points}
\item GitHub commits \emph{- 200 points}
\item GitHub README Markdown page \emph{- 50 points}
\item \LaTeX\ summary report \emph{- 50 points}
\item Project proposal presentation \emph{- 200 points}
\item Project progress presentation \emph{- 200 points}
\item Project final report presentation \emph{- 300 points}
\item Total \emph{- 1200 points}
\end{itemize}

\begin{center}
\begin{tabular}{|c|c||c|c|}
	\hline
	Grade &  Percent Range & Grade  & Percent Range   \\
	\hline
	\hline
	A & 91-100 &  C &  61-65 \\
	\hline
	A-&  89-90&  C-&  56-60  \\
	\hline
	B+& 81-85 &  D+& 51-55 \\
	\hline
	B & 76-80 &  D&  46-50 \\
	\hline
	B- & 71-75 & D- &  41-45 \\
	\hline
	C+&  66-70&  F & 0-40 \\
	\hline
	\hline
\end{tabular}
\end{center}



\section*{Learning Resources}
All learning resources will be supplied at no cost. 
All used software will be free and open-source.\\ 
For complete access to all teaching materials and learning resources, see\\
\url{https://github.com/analyticalworkflows/TeachingMaterials}.

\clearpage

\section*{Course Policies }


\subsection*{Academic Calendar}
All students are subject to the registration and refund deadlines as stated in the Academic Calendar: 
\url{https://registrar.oregonstate.edu/osu-academic-calendar}

\subsection*{Statement Regarding Students with Disabilities}
Accommodations for students with disabilities are determined and approved by Disability Access 
Services (DAS). 
If you, as a student, believe you are eligible for accommodations but have not obtained 
approval please contact DAS immediately at 541-737-4098 or at \url{http://ds.oregonstate.edu}. 
DAS notifies 
students and faculty members of approved academic accommodations and coordinates implementation 
of those accommodations. 
While not required, students and faculty members are encouraged to discuss 
details of the implementation of individual accommodations.

\subsection*{Student Conduct Expectations}
\url{https://beav.es/codeofconduct}

\subsection*{Reach Out for Success}
University students encounter setbacks from time to time. If you encounter difficulties and need 
assistance, it?s important to reach out. Consider discussing the situation with an instructor or academic 
advisor.
 Learn about resources that assist with wellness and academic success at 
\url{oregonstate.edu/ReachOut}.
 If you are in immediate crisis, please contact the Crisis Text Line by texting 
OREGON to 741-741 or call the National Suicide Prevention Lifeline at 1-800-273-TALK (8255)

\subsection*{Student Evaluation of Courses}
During Fall, Winter, and Spring term, the online Student Evaluation of Teaching system opens to 
students the Wednesday of week 8 and closes the Sunday before Finals Week. 
Students will receive notification, instructions and the link through their ONID email. 
They may also log into the system via Online Services. 
Course evaluation results are extremely important and used to help improve courses and the learning 
experience of future students. 
Responses are anonymous (unless a student chooses to ``sign'' their comments, agreeing to relinquish 
anonymity) and unavailable to instructors until after grades have been posted. 
The results of scaled questions and signed comments go to both the instructor and their unit 
head/supervisor. 
Anonymous (unsigned) comments go to the instructor only.


\end{document}
